\documentclass[a4paper,12pt]{article}
\usepackage[utf8]{inputenc}
\usepackage[T2A]{fontenc}
\usepackage[russian]{babel}
\usepackage{amsmath}
\usepackage{amssymb}
\usepackage{array}
\usepackage{booktabs}
\usepackage{geometry}
\usepackage{xcolor}
\usepackage{tabularx}

\geometry{top=2cm, bottom=2cm, left=2.5cm, right=2cm}

\title{\textbf{Домашняя работа 1 по дискретной математике}\\[0.5em]
\large Раскраска графа методом упорядочивания вершин (алг. 7.2.2)}
\author{Кливцов Александр Александрович \\ ISU: 504590, Вариант 29}
\date{}

\begin{document}

\maketitle
\thispagestyle{empty}

% ──────────────────────────────────────────────
\section*{Условие задачи}

Дан граф $G$ с 12 вершинами $e_1, e_2, \ldots, e_{12}$, заданный матрицей смежности~$R$.
Требуется найти правильную раскраску вершин графа, применив \textbf{алгоритм 7.2.2}
(алгоритм, использующий упорядочивание вершин).

% ──────────────────────────────────────────────
\section*{Исходная матрица смежности $R$}

Пустая ячейка означает отсутствие ребра между вершинами.

\[
	R =
	\renewcommand{\arraystretch}{1.1}
	\begin{array}{c|cccccccccccc}
		       & e_1 & e_2 & e_3 & e_4 & e_5 & e_6 & e_7 & e_8 & e_9 & e_{10} & e_{11} & e_{12} \\\hline
		e_1    & 0   & 4   & 4   &     & 2   &     &     & 2   &     &        & 3      &        \\
		e_2    & 4   & 0   &     &     & 4   &     & 2   &     & 5   &        & 2      & 2      \\
		e_3    & 4   &     & 0   &     &     & 4   & 4   & 4   &     & 5      & 2      &        \\
		e_4    &     &     &     & 0   &     & 5   &     &     & 4   &        & 4      &        \\
		e_5    & 2   & 4   &     &     & 0   & 3   &     &     &     &        &        & 4      \\
		e_6    &     &     & 4   & 5   & 3   & 0   &     & 4   &     & 3      & 2      &        \\
		e_7    &     & 2   & 4   &     &     &     & 0   &     &     & 1      & 3      &        \\
		e_8    & 2   &     & 4   &     &     & 4   &     & 0   &     & 1      & 3      &        \\
		e_9    &     & 5   &     & 4   &     &     &     &     & 0   &        &        &        \\
		e_{10} &     &     & 5   &     &     & 3   & 1   & 1   &     & 0      &        &        \\
		e_{11} & 3   & 2   & 2   & 4   &     & 2   & 3   & 3   &     &        & 0      &        \\
		e_{12} &     & 2   &     &     & 4   &     &     &     &     &        &        & 0      \\
	\end{array}
\]

% ──────────────────────────────────────────────
\section*{Список смежности}

Из матрицы выписываем множества смежных вершин:

\begin{center}
	\begin{tabular}{ll}
		$N(e_1) = \{e_2, e_3, e_5, e_8, e_{11}\}$         &
		$N(e_7) = \{e_2, e_3, e_{10}, e_{11}\}$             \\
		$N(e_2) = \{e_1, e_5, e_7, e_9, e_{11}, e_{12}\}$ &
		$N(e_8) = \{e_1, e_3, e_6, e_{10}, e_{11}\}$        \\
		$N(e_3) = \{e_1, e_6, e_7, e_8, e_{10}, e_{11}\}$ &
		$N(e_9) = \{e_2, e_4\}$                             \\
		$N(e_4) = \{e_6, e_9, e_{11}\}$                   &
		$N(e_{10}) = \{e_3, e_6, e_7, e_8\}$                \\
		$N(e_5) = \{e_1, e_2, e_6, e_{12}\}$              &
		$N(e_{11}) = \{e_1, e_2, e_3, e_4, e_6, e_7, e_8\}$ \\
		$N(e_6) = \{e_3, e_4, e_5, e_8, e_{10}, e_{11}\}$ &
		$N(e_{12}) = \{e_2, e_5\}$                          \\
	\end{tabular}
\end{center}

% ──────────────────────────────────────────────
\section*{Выполнение алгоритма 7.2.2}

% ══════════════════════
\subsection*{Итерация $j = 1$}

\textbf{Шаг 2. Подсчёт числа ненулевых элементов $r_i$ (степени вершин):}

\begin{center}
	\begin{tabular}{|c|*{12}{c|}}
		\hline
		Вершина & $e_1$ & $e_2$ & $e_3$ & $e_4$ & $e_5$ & $e_6$ & $e_7$ & $e_8$ & $e_9$ & $e_{10}$ & $e_{11}$   & $e_{12}$ \\
		\hline
		$r_i$   & 5     & 6     & 6     & 3     & 4     & 6     & 4     & 5     & 2     & 4        & \textbf{7} & 2        \\
		\hline
	\end{tabular}
\end{center}

\textbf{Шаг 3. Упорядочивание вершин по невозрастанию $r_i$:}
\[
	e_{11}(7),\; e_2(6),\; e_3(6),\; e_6(6),\; e_1(5),\; e_8(5),\;
	e_5(4),\; e_7(4),\; e_{10}(4),\; e_4(3),\; e_9(2),\; e_{12}(2)
\]

\textbf{Шаг 4. Раскраска в цвет $j = 1$} (просматриваем слева направо):

\begin{enumerate}
	\item $e_{11}$: неокрашена, соседей цвета~1 нет
	      $\;\Rightarrow\;$ \textbf{окрашиваем в цвет~1}.
	\item $e_2$: $e_{11} \in N(e_2)$, $e_{11}$ уже имеет цвет~1 $\;\Rightarrow\;$ пропускаем.
	\item $e_3$: $e_{11} \in N(e_3)$ $\;\Rightarrow\;$ пропускаем.
	\item $e_6$: $e_{11} \in N(e_6)$ $\;\Rightarrow\;$ пропускаем.
	\item $e_1$: $e_{11} \in N(e_1)$ $\;\Rightarrow\;$ пропускаем.
	\item $e_8$: $e_{11} \in N(e_8)$ $\;\Rightarrow\;$ пропускаем.
	\item $e_5$: $N(e_5) = \{e_1, e_2, e_6, e_{12}\}$; ни одна не окрашена в цвет~1
	      $\;\Rightarrow\;$ \textbf{окрашиваем в цвет~1}.
	\item $e_7$: $e_{11} \in N(e_7)$ $\;\Rightarrow\;$ пропускаем.
	\item $e_{10}$: $N(e_{10}) = \{e_3, e_6, e_7, e_8\}$; не смежна с $e_{11}$ и $e_5$
	      $\;\Rightarrow\;$ \textbf{окрашиваем в цвет~1}.
	\item $e_4$: $e_{11} \in N(e_4)$ $\;\Rightarrow\;$ пропускаем.
	\item $e_9$: $N(e_9) = \{e_2, e_4\}$; не смежна с $e_{11}, e_5, e_{10}$
	      $\;\Rightarrow\;$ \textbf{окрашиваем в цвет~1}.
	\item $e_{12}$: $e_5 \in N(e_{12})$, $e_5$ имеет цвет~1 $\;\Rightarrow\;$ пропускаем.
\end{enumerate}

\[
	\boxed{\text{Цвет~1} = \{e_5,\; e_9,\; e_{10},\; e_{11}\}}
\]

Неокрашенные вершины: $e_1, e_2, e_3, e_4, e_6, e_7, e_8, e_{12}$.

% ══════════════════════
\subsection*{Итерация $j = 2$}

\textbf{Шаг 5.} Удаляем из $R$ строки и столбцы $e_5, e_9, e_{10}, e_{11}$. Полагаем $j = 2$.

\textbf{Шаг 2. Пересчёт $r_i$ по подматрице} (вершины $e_1, e_2, e_3, e_4, e_6, e_7, e_8, e_{12}$):

\begin{itemize}
	\item $e_1$: $\{e_2, e_3, e_8\}$ (удалены $e_5, e_{11}$) $\Rightarrow r_1 = 3$
	\item $e_2$: $\{e_1, e_7, e_{12}\}$ (удалены $e_5, e_9, e_{11}$) $\Rightarrow r_2 = 3$
	\item $e_3$: $\{e_1, e_6, e_7, e_8\}$ (удалены $e_{10}, e_{11}$) $\Rightarrow r_3 = 4$
	\item $e_4$: $\{e_6\}$ (удалены $e_9, e_{11}$) $\Rightarrow r_4 = 1$
	\item $e_6$: $\{e_3, e_4, e_8\}$ (удалены $e_5, e_{10}, e_{11}$) $\Rightarrow r_6 = 3$
	\item $e_7$: $\{e_2, e_3\}$ (удалены $e_{10}, e_{11}$) $\Rightarrow r_7 = 2$
	\item $e_8$: $\{e_1, e_3, e_6\}$ (удалены $e_{10}, e_{11}$) $\Rightarrow r_8 = 3$
	\item $e_{12}$: $\{e_2\}$ (удалена $e_5$) $\Rightarrow r_{12} = 1$
\end{itemize}

\textbf{Шаг 3. Упорядочивание:}
\[
	e_3(4),\; e_1(3),\; e_2(3),\; e_6(3),\; e_8(3),\; e_7(2),\; e_4(1),\; e_{12}(1)
\]

\textbf{Шаг 4. Раскраска в цвет $j = 2$:}

\begin{enumerate}
	\item $e_3$: неокрашена, соседей цвета~2 нет $\;\Rightarrow\;$ \textbf{окрашиваем в цвет~2}.
	\item $e_1$: $e_3 \in N(e_1)$, $e_3$ имеет цвет~2 $\;\Rightarrow\;$ пропускаем.
	\item $e_2$: $N(e_2) \cap \{\text{цвет~2}\} = \emptyset$ (соседи: $e_1, e_7, e_{12}$; $e_3 \notin N(e_2)$)
	      $\;\Rightarrow\;$ \textbf{окрашиваем в цвет~2}.
	\item $e_6$: $e_3 \in N(e_6)$ $\;\Rightarrow\;$ пропускаем.
	\item $e_8$: $e_3 \in N(e_8)$ $\;\Rightarrow\;$ пропускаем.
	\item $e_7$: $e_3 \in N(e_7)$ и $e_2 \in N(e_7)$ $\;\Rightarrow\;$ пропускаем.
	\item $e_4$: $N(e_4) = \{e_6\}$; не смежна с $e_3$ или $e_2$
	      $\;\Rightarrow\;$ \textbf{окрашиваем в цвет~2}.
	\item $e_{12}$: $e_2 \in N(e_{12})$, $e_2$ имеет цвет~2 $\;\Rightarrow\;$ пропускаем.
\end{enumerate}

\[
	\boxed{\text{Цвет~2} = \{e_2,\; e_3,\; e_4\}}
\]

Неокрашенные вершины: $e_1, e_6, e_7, e_8, e_{12}$.

% ══════════════════════
\subsection*{Итерация $j = 3$}

\textbf{Шаг 5.} Удаляем из подматрицы строки и столбцы $e_2, e_3, e_4$. Полагаем $j = 3$.

\textbf{Шаг 2. Пересчёт $r_i$} (вершины $e_1, e_6, e_7, e_8, e_{12}$):

\begin{itemize}
	\item $e_1$: $\{e_8\}$ (удалены $e_2, e_3$) $\Rightarrow r_1 = 1$
	\item $e_6$: $\{e_8\}$ (удалены $e_3, e_4$) $\Rightarrow r_6 = 1$
	\item $e_7$: $\{\}$ (удалены $e_2, e_3$) $\Rightarrow r_7 = 0$
	\item $e_8$: $\{e_1, e_6\}$ (удален $e_3$) $\Rightarrow r_8 = 2$
	\item $e_{12}$: $\{\}$ (удалена $e_2$) $\Rightarrow r_{12} = 0$
\end{itemize}

\textbf{Шаг 3. Упорядочивание:}
\[
	e_8(2),\; e_1(1),\; e_6(1),\; e_7(0),\; e_{12}(0)
\]

\textbf{Шаг 4. Раскраска в цвет $j = 3$:}

\begin{enumerate}
	\item $e_8$: неокрашена $\;\Rightarrow\;$ \textbf{окрашиваем в цвет~3}.
	\item $e_1$: $e_8 \in N(e_1)$ $\;\Rightarrow\;$ пропускаем.
	\item $e_6$: $e_8 \in N(e_6)$ $\;\Rightarrow\;$ пропускаем.
	\item $e_7$: в подграфе соседей нет; $e_7 \not\sim e_8$
	      $\;\Rightarrow\;$ \textbf{окрашиваем в цвет~3}.
	\item $e_{12}$: в подграфе соседей нет; $e_{12} \not\sim e_8$, $e_{12} \not\sim e_7$
	      $\;\Rightarrow\;$ \textbf{окрашиваем в цвет~3}.
\end{enumerate}

\[
	\boxed{\text{Цвет~3} = \{e_7,\; e_8,\; e_{12}\}}
\]

Неокрашенные вершины: $e_1, e_6$.

% ══════════════════════
\subsection*{Итерация $j = 4$}

\textbf{Шаг 5.} Удаляем $e_7, e_8, e_{12}$. Полагаем $j = 4$.

\textbf{Шаг 2. Пересчёт $r_i$} (вершины $e_1, e_6$):

$e_1$ и $e_6$ не смежны между собой (в исходной матрице $r_{1,6}$ отсутствует),
поэтому $r_1 = 0$, $r_6 = 0$.

\textbf{Шаг 3. Упорядочивание:}
\[
	e_1(0),\; e_6(0)
\]

\textbf{Шаг 4. Раскраска в цвет $j = 4$:}

\begin{enumerate}
	\item $e_1$: неокрашена $\;\Rightarrow\;$ \textbf{окрашиваем в цвет~4}.
	\item $e_6$: $e_1 \notin N(e_6)$ $\;\Rightarrow\;$ \textbf{окрашиваем в цвет~4}.
\end{enumerate}

\[
	\boxed{\text{Цвет~4} = \{e_1,\; e_6\}}
\]

Все вершины окрашены. \textbf{Задача решена.}

% ──────────────────────────────────────────────
\section*{Итоговая раскраска}

\begin{center}
	\renewcommand{\arraystretch}{1.3}
	\begin{tabular}{|c|l|c|}
		\hline
		\textbf{Цвет}                               & \textbf{Вершины}                 & \textbf{Количество вершин} \\
		\hline
		1                                           & $e_5,\; e_9,\; e_{10},\; e_{11}$ & 4                          \\
		2                                           & $e_2,\; e_3,\; e_4$              & 3                          \\
		3                                           & $e_7,\; e_8,\; e_{12}$           & 3                          \\
		4                                           & $e_1,\; e_6$                     & 2                          \\
		\hline
		\multicolumn{2}{|c|}{\textbf{Всего вершин}} & \textbf{12}                                                   \\
		\hline
	\end{tabular}
\end{center}

\[
	\boxed{\chi(G) = 4}
\]

\textbf{Хроматическое число графа равно 4.}

\end{document}
